\documentclass[UTF8,a4paper]{ctexart}
\usepackage{amsmath,amssymb,amsthm,graphicx,hyperref}
\usepackage{geometry}
\geometry{a4paper,margin=2.5cm}
\title{前言与系列介绍}
\author{《机器人操作器控制：理论与实践》中文译本}
\date{}

\begin{document}
\maketitle

\section{系列介绍}

许多关于控制工程的教科书已经出版，描述了控制系统的新技术，或以新的、更好的方式从数学上阐述现有方法，以解决实际工程师面临的日益复杂的问题。然而，这些书籍很少涉及控制工程的应用方面。这个新系列的目的就是纠正这种情况。

该系列将强调应用问题，而不仅仅是控制工程的数学。它将提供既介绍新技术又介绍成熟技术的教材，同时提供这些方法应用于解决现实世界问题的详细示例。作者将来自学术界和相关应用领域。

在电气、机械（包括航空航天）和化学工程等成熟领域，已经有许多控制技术应用的精彩示例。我们只需环顾当今高度自动化的社会，就能看到制造业中先进机器人技术的应用；空中和地面运输系统中自动控制和导航系统的使用；家用消费品中智能控制系统日益增多的使用；以及对家庭消费者和工业的可靠供水、供气和供电。然而，目前有许多具有挑战性的问题，如果能更广泛地应用控制方法学和控制系统化的应用基础，将会受益匪浅。

本系列推出的书籍将借鉴学术界和应用领域的专业知识，不仅可作为学术界推荐的课程教材，也可作为应用领域从业者的手册。《非线性控制系统》是Dekker控制工程系列中的另一部杰出作品。

\section{前言}

"机器人"一词由捷克剧作家卡雷尔·恰佩克（Karel Capek）在他1920年的戏剧《罗素姆的万能机器人》（Rossum's Universal Robots）中首次引入。捷克语中的"robota"一词简单地意味着"工作"。尽管有如此实用的起源，科幻小说作家和早期好莱坞电影却给了我们一种关于机器人的浪漫概念。这些机器的人形特性似乎给机器人的概念引入了某种人类寻找自身身份的元素。

"自动化"一词于20世纪40年代在福特汽车公司引入，是"自动动机"的缩写。"自动化"这一单一术语将两个概念结合在一起：为机械执行任务而设计的专用机器人机器的概念，以及指挥它们的自动控制系统的概念。

自动控制系统的历史有着深厚的根源。希腊人和阿拉伯人的大多数反馈控制器用于调节水钟以准确报时；这些在14世纪瑞士发明机械钟后被淘汰。自动控制系统直到三百年后的工业革命时期，随着需要先进控制器的复杂机器的出现，才真正兴起；我们特别想到的是风力磨坊和蒸汽机。另一方面，尽管由其他人（如T.Newcomen于1712年）发明，但蒸汽机的荣誉通常归于詹姆斯·瓦特，他在1769年制造了他的发动机，将机械创新与允许自动调节的控制系统结合在一起。也就是说，现代复杂机器除非配备合适的控制系统，否则是没有用的。

瓦特于1788年设计的离心飞球调速器提供了恒速控制器，使蒸汽机能够在工业中高效使用。飞球调速器的运动即使对未经训练的眼睛也是清晰可见的，其原理对许多人来说似乎具有异国情调，体现了新时代的精神。因此，调速器很快在欧洲各地引起了轰动。

主从遥操作机制于20世纪40年代中期在橡树岭和阿贡国家实验室用于放射性物质的远程处理。第一台商用机器人于20世纪50年代末由Unimation公司推出（几乎与1957年的斯普特尼克同时——因此太空时代和机器人时代同时开始）。像飞球调速器一样，机器人操作器的运动对未经训练的眼睛也是显而易见的，因此机器人设备的潜力可以激发想象力。然而，20世纪60年代对工业和非结构化环境中自主机器人自动化的高期望通常未能实现。这是因为今天的机器人学与1712年纽科门工作后不久蒸汽机所处的阶段相同。

机器人学是一个跨学科领域，涉及物理学、机械设计、静力学和动力学、电子学、控制理论、传感器、视觉、信号处理、计算机编程、人工智能（AI）和制造等多种学科。各种专家研究机器人学的各种有限方面，但很少有工程师能够同时面对所有这些领域。这进一步增加了机器人学的浪漫化性质，例如，对于控制理论家来说，他对AI有着堂吉诃德式的幻想。

我们可以将机器人学分为五个主要领域：运动控制、传感器和视觉、规划和协调、人工智能和决策，以及人机接口。没有良好的控制系统，机器人设备是无用的。机器人手臂及其控制系统可以被封装为通用数据抽象；也就是说，机器人加控制器被视为与外部世界交互的单一实体或"代理"。

机器人代理的能力由运动和力施加能力的机械精度、手臂的自由度数量、抓手的可操控程度、传感器以及控制器的复杂性和可靠性决定。机器人手臂的输入仅仅是电机电流和电压，或液压或气动压力；然而，机器人加控制器代理的输入可以是期望的运动轨迹，或期望的施加力。因此，控制系统将机器人在抽象层次上提升了一个层次。

本书旨在提供对串联连杆机器人手臂控制系统的深入研究。它是我们1993年书籍的修订和扩展版本。增加了关于商用机器人操作器和设备、神经网络智能控制以及在真实机器人系统上实现先进控制器的新章节。第1章通过描述现有的机器人及其局限性和能力、传感器和控制器，将本书置于现有商业机器人系统的背景中。

我们希望本书既适合控制工程师也适合机器人学家。因此，附录A提供了机器人运动学和雅可比矩阵的背景知识，第2章提供了控制理论和数学概念的背景知识。本书的初衷是作为研究生第二门机器人课程的教材，但鉴于背景材料，在UTA它被用作电气工程研究生的第一年课程。该课程也被列为本科课程的一部分，本科生很快就掌握了这些内容。

第3章介绍了作为控制设计基础所需的机器人动力学方程。附录C和全书示例中给出了一些常见手臂的动力学。第4章涵盖了计算力矩控制这一重要主题，它提供了重要的见解，同时也在统一框架中汇集了几种经典和现代机器人控制方案。

鲁棒控制和自适应控制分别在第5章和第6章中以并行方式介绍，以突出这两种在不确定性和干扰面前的控制方法的相似之处和不同之处。第7章讨论了一些先进技术，包括学习控制和具有柔性关节耦合的手臂。

基于生物系统的现代智能控制技术已经解决了复杂系统控制中的许多问题，包括未知的非参数化动力学和未知干扰、间隙、摩擦和死区。因此，我们在第8章增加了一章关于神经网络控制系统的内容。机器人只有在其与外部环境接触时才有用，因此第9章讨论了力控制问题。

成功控制器设计验证的关键是计算机仿真。因此，我们在全书中举例说明了受控非线性系统的计算机仿真过程。附录B提供了仿真软件。MATLAB等商业软件包使机器人控制系统的仿真变得非常容易。

设计了机器人控制系统后，有必要实现它；鉴于今天的微处理器和数字信号处理器，从计算机仿真到实现只是一小步，因为仿真所需的控制器子程序（包含在本书中）与实际手臂实现所需的微处理器中的子程序几乎相同。事实上，第10章展示了在真实机器人系统上实现本书开发的先进控制器的技术。

所有基本信息和控制设计算法都以表格形式显示在书中。这连同书前的示例列表和表格列表，为学生、学者或实际工程师提供了方便的参考。

我们感谢Milagro Design的Wei Cheng负责本书的LaTeX排版和图表制作，以及将第一版的内容扫描成电子格式。

\begin{flushright}
F.L. Lewis，德克萨斯州阿灵顿\\
D.M. Dawson，南卡罗来纳州克莱姆森\\
C.T. Abdallah，新墨西哥州阿尔伯克基
\end{flushright}

\section{目录结构}

\begin{tabular}{ll}
系列介绍 & v \\
前言 & vii \\
第1章　商用机器人操作器 & 1 \\
\quad 1.1 引言 & 1 \\
\quad 1.2 商用机器人构型与类型 & 3 \\
\quad 1.3 商用机器人控制器 & 10 \\
\quad 1.4 传感器 & 12 \\
第2章　控制理论导论 & 21 \\
\quad 2.1 引言 & 21 \\
\quad 2.2 线性状态变量系统 & 22 \\
\quad 2.3 非线性状态变量系统 & 31 \\
\quad 2.4 非线性系统与平衡点 & 36 \\
\quad 2.5 向量空间、范数与内积 & 39 \\
\quad 2.6 稳定性理论 & 51 \\
\quad 2.7 李雅普诺夫稳定性定理 & 67 \\
\quad 2.8 输入/输出稳定性 & 80 \\
\quad 2.9 高级稳定性结果 & 82 \\
\quad 2.10 有用的定理与引理 & 88 \\
\quad 2.11 线性控制器设计 & 93 \\
\quad 2.12 总结与注记 & 101 \\
\end{tabular}

\newpage

\begin{tabular}{ll}
第3章　机器人动力学 & 107 \\
\quad 3.1 引言 & 107 \\
\quad 3.2 拉格朗日-欧拉动力学 & 108 \\
\quad 3.3 机器人方程的结构与性质 & 125 \\
\quad 3.4 状态变量表示与反馈线性化 & 142 \\
\quad 3.5 笛卡尔动力学与其他动力学 & 148 \\
\quad 3.6 执行器动力学 & 152 \\
\quad 3.7 总结 & 161 \\
第4章　计算力矩控制 & 169 \\
\quad 4.1 引言 & 169 \\
\quad 4.2 路径生成 & 170 \\
\quad 4.3 机器人系统的计算机仿真 & 181 \\
\quad 4.4 计算力矩控制 & 185 \\
\quad 4.5 数字机器人控制 & 222 \\
\quad 4.6 最优外环设计 & 243 \\
\quad 4.7 笛卡尔控制 & 248 \\
\quad 4.8 总结 & 251 \\
第5章　机器人操作器的鲁棒控制 & 263 \\
\quad 5.1 引言 & 263 \\
\quad 5.2 反馈线性化控制器 & 265 \\
\quad 5.3 非线性控制器 & 293 \\
\quad 5.4 动力学重新设计 & 316 \\
\quad 5.5 总结 & 320 \\
\end{tabular}

\newpage

\begin{tabular}{ll}
第6章　机器人操作器的自适应控制 & 329 \\
\quad 6.1 引言 & 329 \\
\quad 6.2 计算力矩方法的自适应控制 & 330 \\
\quad 6.3 惯量相关方法的自适应控制 & 341 \\
\quad 6.4 基于无源性的自适应控制器 & 349 \\
\quad 6.5 持续激励 & 357 \\
\quad 6.6 复合自适应控制器 & 361 \\
\quad 6.7 自适应控制器的鲁棒性 & 371 \\
\quad 6.8 总结 & 377 \\
第7章　先进控制技术 & 383 \\
\quad 7.1 引言 & 383 \\
\quad 7.2 减少在线计算的机器人控制器 & 384 \\
\quad 7.3 自适应鲁棒控制 & 399 \\
\quad 7.4 执行器动力学补偿 & 407 \\
\quad 7.5 总结 & 426 \\
第8章　机器人的神经网络控制 & 431 \\
\quad 8.1 引言 & 431 \\
\quad 8.2 神经网络背景 & 433 \\
\quad 8.3 使用静态神经网络的跟踪控制 & 440 \\
\quad 8.4 线性参数神经网络的调参算法 & 445 \\
\quad 8.5 非线性参数神经网络的调参算法 & 449 \\
\quad 参考文献 & 459 \\
\end{tabular}

\newpage

\begin{tabular}{ll}
第9章　力控制 & 463 \\
\quad 9.1 引言 & 463 \\
\quad 9.2 刚度控制 & 464 \\
\quad 9.3 混合位置/力控制 & 478 \\
\quad 9.4 混合阻抗控制 & 489 \\
\quad 9.5 降维状态位置/力控制 & 501 \\
\quad 9.6 总结 & 510 \\
第10章　机器人控制实现与软件 & 517 \\
\quad 10.1 引言 & 518 \\
\quad 10.2 工具与技术 & 520 \\
\quad 10.3 机器人平台设计 & 523 \\
\quad 10.4 机器人平台操作 & 543 \\
\quad 10.5 编程示例 & 548 \\
\quad 10.6 总结 & 550 \\
附录A　机器人运动学与雅可比矩阵回顾 & 555 \\
附录B　控制器仿真软件 & 591 \\
附录C　一些常见机器人手臂的动力学 & 599 \\
\end{tabular}

\end{document}
